% ============================================================================
% Fichier  : partie3/P03_C11_legislation_mondiale.tex
% Titre    : Législation Mondiale et Régionale
% Auteur   : MINKA MI NGUIDJOÏ Thierry Emmanuel
% Version  : 1.0 -- Juin 2026
% Statut   : RÉDIGÉ
% Sources  : 13_legislation_mondiale_regionale.tex (squelette FRE/SCA/CFAA/
%            RGPD/Budapest/Malabo/CEDEAO/SADC/EAC)
%            M4_Cadre_Normatif_International_PNC.docx (CLOUD Act, Protocole II,
%            comparatif, coopération)
% Positionnement : Ch. 10 = cadre normatif global (Budapest, Malabo, INTERPOL).
%                  Ce chapitre = impact forensique du droit de CHAQUE grande
%                  juridiction sur les investigations camerounaises.
%                  Question directrice : quand les OPJ requis-itionnent un
%                  fournisseur ou coopèrent avec un partenaire, quel droit
%                  local s'applique chez eux, et qu'est-ce que ça change ?
% ============================================================================

\chapter{Législation Mondiale et Régionale}
\label{chap:leg-mond}

\epigraph{%
    « Cybercrime knows no borders, but our laws still do.
    International cooperation is not just desirable --- it's imperative. »%
}{--- Susan W. Brenner}

\niveauChapitre{Fondamental}{%
    Lecture des Chapitres~\ref{chap:norm-glob} (Cadre normatif global)
    recommandée~: ce chapitre examine les droits nationaux que le
    Chapitre~10 ne couvre pas directement.%
}

\begin{boxobjectifs}
\begin{itemize}
    \item Comprendre le droit américain applicable aux investigations
          impliquant des fournisseurs US~: FRE, SCA, CFAA, CLOUD~Act.
    \item Identifier les tensions entre RGPD européen et investigation
          forensique, et comprendre comment les résoudre en pratique.
    \item Maîtriser eIDAS et NIS~2~: leur impact sur la valeur probatoire
          des preuves numériques et la sécurité des systèmes.
    \item Connaître les cadres régionaux africains~: CEDEAO, SADC, EAC
          et leur utilité opérationnelle pour les OPJ camerounais.
    \item Appliquer une grille de lecture comparative pour choisir
          la stratégie juridique optimale dans une affaire transfrontalière.
\end{itemize}
\end{boxobjectifs}

\bigskip

% ============================================================================
\section*{Question directrice de ce chapitre}
% ============================================================================

Quand un OPJ camerounais réquisitionne Google pour les emails d'un suspect,
c'est le droit américain qui s'applique chez Google, pas le droit camerounais.
Quand des preuves proviennent d'un serveur hébergé en Allemagne, c'est le RGPD
qui contraint le fournisseur. Quand un partenaire sénégalais coopère sur une
affaire transfrontalière, c'est leur propre cadre législatif qui détermine ce
qu'ils peuvent partager.

Ce chapitre répond à une question pratique~: \emph{quel droit local s'applique
chez quel partenaire, et qu'est-ce que cela change concrètement pour mes
demandes et mes preuves~?}

\separateur

% ============================================================================
\section{Droit Américain~: le Cadre des Grands Fournisseurs}
\label{sec:droit-usa}
% ============================================================================

Les plus grands fournisseurs de services numériques mondiaux --- Google,
Meta, Microsoft, Amazon, Apple --- sont des entreprises américaines soumises
au droit fédéral des États-Unis. Comprendre ce droit n'est pas une option
académique~: c'est la condition pour formuler des demandes qui aboutissent.

\subsection{Federal Rules of Evidence (FRE)~: admissibilité aux États-Unis}

Les \termecle{Federal Rules of Evidence}\index{Federal Rules of Evidence}
(FRE) régissent l'admissibilité des preuves devant les juridictions fédérales
américaines. Deux règles sont particulièrement importantes pour la forensique
numérique internationale~:

\begin{boxjuris}
\textbf{FRE Rule~901 --- Authentication~:}
\begin{quote}
To satisfy the requirement of authenticating or identifying an item of
evidence, the proponent must produce evidence sufficient to support a finding
that the item is what the proponent claims it is.
\end{quote}
\textbf{Application numérique~:} les hash SHA-256 sont acceptés comme preuve
d'authenticité~; la chaîne de custody doit être documentée~; le témoignage
d'un expert est souvent requis pour expliquer les méthodes forensiques.
MD5 est considéré \emph{deprecated} pour les nouveaux dossiers.
\end{boxjuris}

\begin{boxjuris}
\textbf{FRE Rule~803(6) --- Business Records Exception~:}
Les enregistrements de routine d'une entreprise (logs système, emails
conservés, transactions financières) sont admissibles même sans témoignage
direct si~: (1)~ils ont été créés dans le cours normal des affaires~;
(2)~ils ont été conservés selon les procédures habituelles de l'entreprise~;
(3)~un représentant qualifié en atteste.\\[0.3ex]
\textbf{Impact forensique~:} les réponses aux réquisitions LEA de Meta ou
Google constituent des ``business records'' --- elles sont directement
admissibles selon cette règle devant les tribunaux américains, et leur
valeur probatoire est reconnue dans la plupart des juridictions francophones
par réciprocité.
\end{boxjuris}

\subsection{Stored Communications Act (SCA)~: protection et accès}

Le \termecle{Stored Communications Act}\index{SCA} (SCA, 18~U.S.C.
\S\S~2701--2712) protège les communications électroniques stockées par des
tiers. Il régit ce que les fournisseurs américains \emph{peuvent} et
\emph{doivent} communiquer aux forces de l'ordre.

\begin{tableaucours}{p{3.5cm}p{3.5cm}p{4.5cm}}{%
    \tableauHeader{Type de données} &
    \tableauHeader{Instrument requis} &
    \tableauHeader{Impact pour les OPJ camerounais}
}
Données d'abonné de base (nom, adresse, facturation) &
    Administrative subpoena &
    Accessibles via portail LEA sur simple réquisition~;
    délai~: 7--21~jours \\
Adresses IP de connexion, logs de session &
    Court order (18~U.S.C.~\S~2703(d)) &
    Portail LEA~; urgence~: 24--72h~; standard~: 7--30j \\
Contenu des communications (emails, messages) &
    Search warrant (mandat de perquisition) &
    \textbf{MLAT USA obligatoire}~; délai~: 6--18~mois~;
    Protocole~II Budapest peut réduire ce délai (si ratifié) \\
\end{tableaucours}
\vspace{-0.8em}
\begin{center}{\small\itshape Tableau~11.1 --- SCA~: instruments selon le type de données}\end{center}

\begin{boiteRemarque}{Pourquoi les portails LEA ne donnent pas le contenu des emails}
Le SCA interdit aux fournisseurs américains de communiquer le contenu des
communications sans warrant (mandat de perquisition), même à des autorités
étrangères. Ce n'est pas un refus de coopérer --- c'est une obligation
légale. La seule voie légale pour obtenir le contenu est le MLAT ou le
Protocole~II Budapest (si l'État requis l'a ratifié). Reformuler la demande
en ciblant les métadonnées (admissibles sans warrant) est souvent la
stratégie plus rapide et suffisante.
\end{boiteRemarque}

\subsection{Computer Fraud and Abuse Act (CFAA)~: base des poursuites US}

Le \termecle{CFAA}\index{CFAA} (18~U.S.C.~\S~1030) définit les infractions
informatiques fédérales américaines. Bien que principalement applicable aux
faits commis sur le sol américain, il importe à deux titres~:

\begin{enumerate}
    \item \textbf{Base de coopération}~: quand l'OPJ camerounais demande
          une entraide aux États-Unis, la demande doit démontrer que les
          faits constituent une infraction dans les deux pays
          (\textit{double incrimination}). Montrer que les faits tombent
          sous le CFAA \emph{et} sous la Loi~2010/012 renforce la demande.
    \item \textbf{Compétence extraterritoriale}~: le CFAA s'applique quand
          une infrastructure américaine est visée ou utilisée --- même si
          l'auteur et la victime sont à l'étranger. Un attaquant nigérian
          ciblant une banque camerounaise via des serveurs Amazon peut être
          poursuivi aux États-Unis.
\end{enumerate}

\begin{tableaucours}{p{3.5cm}p{3.0cm}p{5.0cm}}{%
    \tableauHeader{Infraction CFAA} &
    \tableauHeader{Peine max.} &
    \tableauHeader{Équivalent Loi~2010/012}
}
Accès non autorisé à un ordinateur protégé &
    1--10~ans &
    Art.~55 (accès illicite) \\
Accès avec obtention d'information &
    1--10~ans &
    Art.~55 + Art.~78 si préjudice \\
Dommages intentionnels (ransomware) &
    10~ans (1\textsuperscript{re} offense) &
    Art.~63 (atteinte aux données) \\
Fraude informatique &
    5--20~ans &
    Art.~78 (fraude informatique) \\
\end{tableaucours}
\vspace{-0.8em}
\begin{center}{\small\itshape Tableau~11.2 --- CFAA et équivalents Loi~2010/012}\end{center}

\subsection{CLOUD Act~: accès transfrontalier aux données}

Le \termecle{CLOUD Act}\index{CLOUD Act} (\textit{Clarifying Lawful Overseas
Use of Data Act}, 2018) est la loi la plus controversée dans le domaine de
la coopération numérique internationale. Elle détermine les obligations
des entreprises américaines vis-à-vis des demandes étrangères.

\begin{boxCRO}
\textbf{Analyse CRO --- CLOUD Act}

$\mathcal{R}$ (Fiabilité)~: le CLOUD~Act permet aux autorités américaines
d'accéder aux données stockées par des entreprises US \emph{n'importe où
dans le monde}, indépendamment du pays de stockage. Pour la forensique~:
les preuves obtenues via CLOUD~Act ont une valeur probatoire élevée car
elles sont certifiées par l'entreprise américaine.

$\mathcal{O}$ (Opposabilité)~: les données obtenues via CLOUD~Act sont
admissibles devant les tribunaux américains selon FRE~803(6), et
généralement reconnues par réciprocité dans les pays partenaires.

$\mathcal{C}$ (Confidentialité)~: \textbf{tension majeure.} Le CLOUD~Act
signifie que les données d'un ressortissant camerounais hébergées chez
Google aux États-Unis peuvent être accessibles aux autorités américaines
sans que le Cameroun en soit informé. La souveraineté numérique
camerounaise est directement engagée.

\cro{$\downarrow\downarrow$ tension souveraineté}{$\uparrow$ pour preuves obtenues}{$\uparrow\uparrow$ recov. us-based}

\textbf{Implication pratique~:} pour les OPJ camerounais~--- le CLOUD~Act
joue en leur faveur~: les partenaires américains peuvent obtenir des données
chez des fournisseurs US plus facilement qu'un MLAT classique, et les
partager via accords bilatéraux.
\end{boxCRO}

% ============================================================================
\section{Droit Européen~: RGPD, eIDAS, NIS~2 et e-Evidence}
\label{sec:droit-europe}
% ============================================================================

L'Union Européenne a construit le cadre normatif numérique le plus dense
au monde. Pour un OPJ camerounais coopérant avec des partenaires européens
ou réquisitionnant des fournisseurs basés en Europe, quatre textes sont
directement pertinents.

\subsection{RGPD~: tensions avec l'investigation forensique}

Le \termecle{RGPD}\index{RGPD} (Règlement~(UE)~2016/679) est le texte de
référence mondial sur la protection des données personnelles. Il crée des
tensions structurelles avec la forensique numérique~:

\begin{tableaucours}{p{3.5cm}p{3.5cm}p{4.5cm}}{%
    \tableauHeader{Droit RGPD} &
    \tableauHeader{Tension forensique} &
    \tableauHeader{Résolution}
}
Droit à l'effacement (art.~17) &
    Un suspect peut demander la suppression de ses données chez un
    fournisseur européen --- y compris des preuves potentielles &
    Art.~17(3)(e)~: exception pour \textit{constatation, exercice ou
    défense de droits en justice}~; conservation possible sur injonction \\
Minimisation des données (art.~5) &
    La collecte exhaustive de toutes les données d'un suspect viole le
    principe de minimisation &
    L'investigation judiciaire constitue une base légale distincte
    (art.~6(1)(e)~: mission d'intérêt public) \\
Notification de violation (art.~33) &
    Un fournisseur victime d'une fuite doit notifier sous 72h~---
    ce qui peut alerter le suspect visé &
    Les autorités judiciaires peuvent demander le report de la
    notification (art.~23) si cela entrave l'enquête \\
\end{tableaucours}
\vspace{-0.8em}
\begin{center}{\small\itshape Tableau~11.3 --- Tensions RGPD / investigation forensique}\end{center}

\begin{boxjuris}
\textbf{RGPD Article~23 --- Limitations~:}
\begin{quote}
Le droit de l'Union ou le droit d'un État membre \textit{[\ldots]} peut
limiter la portée des obligations \textit{[\ldots]} lorsqu'une telle
limitation respecte l'essence des libertés et droits fondamentaux et qu'elle
constitue une mesure nécessaire et proportionnée dans une société
démocratique pour garantir~: \textit{(d)~la prévention et la détection
d'infractions pénales, ainsi que les enquêtes et les poursuites en la
matière ou l'exécution de sanctions pénales}.
\end{quote}
\textbf{Implication opérationnelle~:} l'article~23 est la base légale
qui permet aux partenaires européens de transmettre des données au
Cameroun malgré le RGPD, dans le cadre d'une coopération judiciaire.
Citer explicitement cet article dans toute demande d'entraide adressée
à un État européen.
\end{boxjuris}

\subsection{eIDAS~: valeur probatoire des signatures et horodatages électroniques}

Le règlement \termecle{eIDAS}\index{eIDAS} (Règlement~(UE)~910/2014,
révisé 2024) établit un cadre pour la confiance numérique en Europe.
Son impact forensique porte sur la valeur probatoire des documents
électroniques provenant d'Europe.

\begin{tableaucours}{p{3.2cm}p{5.5cm}p{3.8cm}}{%
    \tableauHeader{Niveau} &
    \tableauHeader{Description} &
    \tableauHeader{Valeur probatoire}
}
\textbf{Signature simple} &
    Toute donnée électronique liée à un signataire &
    Faible~; contestable \\
\textbf{Signature avancée} &
    Identifie le signataire de manière unique~; liée aux données
    de façon détectable en cas de modification &
    Moyenne~; présomption simple d'authenticité \\
\textbf{Signature qualifiée} &
    Certificat qualifié + dispositif sécurisé certifié QSCD~;
    équivalent légal d'une signature manuscrite &
    \textbf{Maximale}~; présomption irréfragable en droit européen \\
\textbf{Horodatage qualifié} &
    Preuve d'existence d'un document à un instant précis~;
    délivré par une autorité de confiance certifiée &
    \textbf{Force probante maximale} pour la date d'existence \\
\end{tableaucours}
\vspace{-0.8em}
\begin{center}{\small\itshape Tableau~11.4 --- Niveaux de signature et valeur probatoire (eIDAS)}\end{center}

\begin{boiteRemarque}{eIDAS et forensique camerounaise}
Quand un partenaire européen transmet des preuves numériques signées avec
une signature qualifiée eIDAS, ces preuves ont une valeur probatoire
maximale en droit européen. Dans un tribunal camerounais, la \loicam{2024/017}
(art.~6~: principe d'intégrité) et les principes généraux du droit de
la preuve permettent de faire reconnaître cette valeur si l'expert
forensique explique le système de certification. Documenter systématiquement
le niveau eIDAS des documents reçus dans le rapport d'analyse.
\end{boiteRemarque}

\subsection{NIS~2~: implications pour la forensique d'incident}

La Directive \termecle{NIS~2}\index{NIS 2} (Directive~(UE)~2022/2555)
impose des obligations de cybersécurité aux opérateurs de secteurs
essentiels dans l'UE. Son impact sur la forensique~:

\begin{itemize}
    \item \textbf{Obligation de notification d'incident}~: les entités
          essentielles doivent notifier l'incident dans les 24h à l'autorité
          nationale compétente, et dans les 72h un rapport initial complet.
          Ces rapports d'incident constituent des preuves documentaires
          exploitables en forensique.
    \item \textbf{Logs conservés}~: NIS~2 impose la conservation de logs
          de sécurité pendant une durée suffisante pour permettre les
          investigations post-incident. Pour l'OPJ coopérant avec un
          partenaire UE~: ces logs existent, ils sont conservés, ils peuvent
          être requis.
    \item \textbf{Incident Response Plans}~: les entités NIS~2 sont tenues
          de maintenir des procédures de réponse aux incidents documentées.
          La conformité NIS~2 d'un partenaire est une garantie de qualité
          probatoire de ses données.
\end{itemize}

\subsection{Règlement e-Evidence~: la révolution en cours}

Le Règlement~(UE)~2023/1543 (\termecle{e-Evidence}), applicable depuis
2026, crée un mécanisme direct pour l'obtention de preuves électroniques
entre les États membres de l'UE~:

\begin{itemize}
    \item \textbf{European Production Order}~: les autorités judiciaires
          d'un État membre peuvent directement requérir un fournisseur
          établi dans un autre État membre~; délai cible~: 10~jours
          (6h en urgence).
    \item \textbf{Impact pour le Cameroun}~: les délais de coopération
          intra-UE se raccourcissent drastiquement~; les partenaires
          européens du Cameroun peuvent obtenir des données pan-européennes
          bien plus vite, et les partager ensuite dans le cadre d'un MLAT.
\end{itemize}

% ============================================================================
\section{Droit Africain Régional~: CEDEAO, SADC et EAC}
\label{sec:droit-afrique}
% ============================================================================

\subsection{CEDEAO~: le cadre d'Afrique de l'Ouest}

La \termecle{CEDEAO}\index{CEDEAO} (\textit{Communauté Économique des États
de l'Afrique de l'Ouest}) a adopté deux textes harmonisant le droit numérique
de ses 15~membres~:

\begin{tableaucours}{p{4.0cm}p{7.5cm}}{%
    \tableauHeader{Texte} &
    \tableauHeader{Contenu et pertinence forensique}
}
\textbf{Directive C/DIR/1/08/11} sur la cybercriminalité &
    Harmonise les infractions cyberpénales~: accès illicite,
    interception, atteinte à l'intégrité. Infractions similaires
    à Budapest. \textbf{Pertinence~:} base pour demandes de coopération
    vers les membres CEDEAO (Sénégal, Ghana, Nigeria~: membres). \\
\textbf{Acte additionnel A/SA.2/01/10} sur les données personnelles &
    Définit les principes de protection des données pour les États
    membres. Influence les réponses aux réquisitions de données de
    citoyens des États membres. \\
\end{tableaucours}
\vspace{-0.8em}
\begin{center}{\small\itshape Tableau~11.5 --- Textes CEDEAO pertinents pour la forensique}\end{center}

\begin{boiteRemarque}{CEDEAO et coopération avec le Nigeria et le Sénégal}
Le Nigeria et le Sénégal (deux États membres CEDEAO ayant également
ratifié Budapest) constituent les partenaires les plus sollicités par
les OPJ camerounais dans les affaires transfrontalières. La CEDEAO ajoute
une couche de coopération régionale~: des demandes peuvent être acheminées
via le mécanisme CEDEAO \emph{en plus} du canal INTERPOL, augmentant les
chances d'obtenir une réponse rapide.
\end{boiteRemarque}

\subsection{SADC~: le modèle d'Afrique Australe}

La \termecle{SADC}\index{SADC} (\textit{Southern African Development
Community}) a développé une \textit{Model Law on Computer Crime and
Cybercrime} servant de référence pour l'harmonisation législative de ses
16~membres (Zimbabwe, Zambie, Mozambique, Afrique du Sud~\ldots).

Points clés de la \textit{Model Law}~:
\begin{itemize}
    \item Criminalisation cohérente avec Budapest des accès illicites,
          interceptions, atteintes à l'intégrité.
    \item Dispositions spécifiques sur la cyberfraude et les atteintes
          aux données financières~: pertinent pour les fraudes bancaires
          transfrontalières Cameroun-Afrique du Sud.
    \item L'Afrique du Sud (\textit{Cybercrimes Act} 2020) est le système
          le plus abouti de la région~: laboratoire forensique national,
          unité cybercrime SAPS, coopération INTERPOL établie.
\end{itemize}

\subsection{EAC~: le cadre d'Afrique de l'Est}

La \termecle{EAC}\index{EAC} (\textit{East African Community}) a développé
un \textit{Framework for Cyberlaws} couvrant les membres~: Kenya, Tanzanie,
Ouganda, Rwanda, Burundi, Sud-Soudan.

Pour les OPJ camerounais, les deux points de contact pertinents sont~:
\begin{itemize}
    \item \textbf{Rwanda}~: l'un des cadres les plus avancés d'Afrique
          sub-saharienne~; \textit{Computer Crime Law} 2020, Autorité
          nationale de cybersécurité (NCSA) opérationnelle.
    \item \textbf{Kenya}~: \textit{Computer Misuse and Cybercrimes Act}
          2018~; coopération INTERPOL active~; hub financier
          (M-Pesa)~: réquisitions sur les fraudes mobile money
          régionales aboutissent bien.
\end{itemize}

% ============================================================================
\section{Grille de Lecture Comparative~: Choisir sa Stratégie Juridique}
\label{sec:grille-comparative}
% ============================================================================

\subsection{Tableau synthétique~: droit local et implications opérationnelles}

\begin{tableaucours}{p{2.5cm}p{2.8cm}p{3.0cm}p{3.5cm}}{%
    \tableauHeader{Juridiction} &
    \tableauHeader{Texte-clé} &
    \tableauHeader{Spécificité forensique} &
    \tableauHeader{Canal recommandé}
}
\textbf{USA} & SCA + CLOUD~Act &
    Contenu emails~: mandat obligatoire. Métadonnées~: portail LEA. &
    Portail LEA (métadonnées)~; MLAT USA (contenu)~;
    Protocole~II Budapest si disponible \\
\textbf{UE} & RGPD + e-Evidence &
    Art.~23 RGPD~: exception judiciaire. e-Evidence~: 10~jours. &
    MLAT bilatéral~; e-Evidence (intra-UE)~;
    citer art.~23 RGPD explicitement \\
\textbf{Afrique West} & CEDEAO Directive &
    Harmonisation avec Budapest. Sénégal/Ghana/Nigeria~: Budapest. &
    INTERPOL réseau 24/7~; canal CEDEAO en complément \\
\textbf{Afrique Sud} & Malabo~; SADC Model Law &
    Rwanda/Kenya~: lois modernes. AfSud~: coopération établie. &
    INTERPOL~; AFRIPOL~; accords bilatéraux \\
\textbf{Russie / Chine} & (pas Budapest, pas Malabo) &
    Coopération très limitée. Chercher les nœuds intermédiaires. &
    Identifier les segments du chemin dans des États coopérants~;
    OSINT sur les nœuds accessibles \\
\end{tableaucours}
\vspace{-0.8em}
\begin{center}{\small\itshape Tableau~11.6 --- Grille comparative~: droit local et stratégie}\end{center}

\subsection{Processus de qualification d'une affaire transfrontalière}

Face à une affaire impliquant un élément étranger, la démarche en
cinq étapes~:

\begin{enumerate}
    \item \textbf{Identifier les éléments étrangers}~: où est le serveur~?
          Où est le suspect~? Où est hébergé le fournisseur~? Où ont transité
          les fonds~?
    \item \textbf{Cartographier les droits locaux applicables}~: pour chaque
          État impliqué, quel texte encadre la coopération~? Budapest,
          Malabo, CEDEAO, accord bilatéral, ou rien~?
    \item \textbf{Qualifier l'infraction en double incrimination}~: montrer
          que les faits constituent une infraction dans les deux pays. Sans
          double incrimination, la coopération peut être refusée.
    \item \textbf{Prioriser~: données volatiles d'abord}~: activer
          immédiatement la conservation (art.~29 Budapest, portail LEA).
          Le droit local du partenaire ne change rien à cette priorité.
    \item \textbf{Choisir le canal adapté}~: MLAT formel si contenu
          requis~; portail LEA si métadonnées suffisent~; INTERPOL 24/7
          si urgence.
\end{enumerate}

\begin{boxscenario}{Affaire BEC Cameroun-UE-USA --- stratégie juridique}
\textbf{Faits~:} SONATRAV SA, Douala, perd 85~millions FCFA dans une
attaque BEC. L'email frauduleux a été envoyé depuis une adresse Gmail
(\texttt{@gmail.com}), via un VPN hébergé en Allemagne. Les fonds ont
transité vers un compte Revolut (UE) puis vers Binance (Caïmans).

\textbf{Cartographie juridique~:}
\begin{itemize}
    \item Gmail~$\to$ \textbf{Google USA}~: SCA~; métadonnées via portail
          LEA~; contenu MLAT USA~; base~: CFAA~\S~1030 (fraude via
          ordinateur protégé) + Loi~2010/012 art.~78.
    \item VPN Allemagne~$\to$ \textbf{droit allemand + RGPD}~: MLAT
          franco-camerounais indirect (via UE)~; citer art.~23 RGPD~;
          e-Evidence~\S~si partenaire UE impliqué.
    \item Revolut (Lituanie/UE)~$\to$ \textbf{RGPD + e-Evidence}~:
          European Production Order~; délai~10~jours.
    \item Binance~$\to$ identifier la juridiction réelle~; analyse
          blockchain Blockchair~; portail LEA Binance si disponible.
\end{itemize}

\textbf{Séquence d'actions~:}
\begin{enumerate}
    \item Conservation urgente simultanée~: portail LEA Google (urgence)
          + BCN INTERPOL pour VPN Allemagne (art.~29 Budapest).
    \item Analyse blockchain Binance~: tracer les fonds (OSINT légal).
    \item MLAT USA pour contenu Gmail dans les 60~jours.
    \item Demande e-Evidence à Revolut Lituanie.
\end{enumerate}
\end{boxscenario}

\separateur

\begin{boxresume}
\textbf{Ce qu'il faut retenir de ce chapitre~:}
\begin{itemize}
    \item \textbf{Droit USA}~: SCA impose un \emph{mandat} pour le contenu
          des communications~; les métadonnées s'obtiennent via portail LEA.
          CLOUD~Act~: les entreprises US peuvent livrer des données stockées
          n'importe où. CFAA~: base pour la double incrimination.

    \item \textbf{RGPD européen}~: tension structurelle avec la forensique,
          résolue par l'art.~23 (exception judiciaire) et l'art.~6(1)(e)
          (mission d'intérêt public). Citer explicitement ces articles
          dans toute demande à un partenaire européen.

    \item \textbf{eIDAS}~: trois niveaux de signature~; la signature
          qualifiée a valeur légale maximale~; l'horodatage qualifié
          prouve l'existence à un instant~; documenter le niveau eIDAS
          des preuves reçues d'Europe.

    \item \textbf{NIS~2}~: impose la conservation de logs et la
          notification d'incidents~; les entités conformes constituent
          des sources de preuves documentaires de qualité.

    \item \textbf{Afrique régionale}~: CEDEAO renforce les demandes vers
          Sénégal/Ghana/Nigeria~; SADC couvre l'Afrique australe (Afrique
          du Sud~: meilleur partenaire forensique)~; EAC couvre
          Rwanda/Kenya (lois modernes, coopération active).

    \item \textbf{Double incrimination}~: condition nécessaire dans la
          plupart des coopérations~; montrer que les faits constituent une
          infraction chez les deux partenaires. Loi~2010/012 +
          CFAA/RGPD-infraction/Budapest = base solide.

    \item \textbf{Grille de décision en 5 étapes}~: identifier les
          éléments étrangers~$\to$ cartographier les droits~$\to$ double
          incrimination~$\to$ conservation d'abord~$\to$ canal adapté.
\end{itemize}
\end{boxresume}

\bigskip

\begin{boxexo}[title={\ding{45}\quad Exercice~11.1 --- Identification du cadre légal \niveaubase}]
Pour chacune des situations suivantes, identifiez~: (A)~le ou les textes
légaux applicables chez le partenaire étranger, (B)~l'instrument requis
pour obtenir les données demandées, (C)~le délai réaliste~:
\begin{enumerate}[(a)]
    \item Vous souhaitez obtenir les données d'abonné d'un compte Gmail
          suspect (\texttt{@gmail.com}) utilisé dans une fraude MoMo.
    \item Vous devez obtenir le contenu des emails Outlook d'un suspect
          hébergés par Microsoft USA dans le cadre d'une escroquerie.
    \item Un serveur Apache hébergé en France a servi à déposer un malware.
          Vous demandez les logs d'accès au fournisseur d'hébergement.
    \item Un compte Facebook utilisé pour harcèlement identifie son
          propriétaire au Sénégal. Vous demandez les données d'abonné.
\end{enumerate}
\end{boxexo}

\begin{boxexo}[title={\ding{45}\quad Exercice~11.2 --- Tensions RGPD \niveaumoyen}]
Un fournisseur d'hébergement allemand vous informe qu'il ne peut pas
répondre à votre réquisition en raison du RGPD et du principe de
minimisation des données~:
\begin{enumerate}[(a)]
    \item Quel article du RGPD invoqueriez-vous pour contester ce refus~?
          Rédigez le paragraphe de réponse à inclure dans votre demande.
    \item Quelle condition doit être remplie pour que l'exception judiciaire
          de l'art.~23 RGPD soit valablement invoquée~?
    \item Si le fournisseur maintient son refus, par quel canal alternatif
          pouvez-vous obtenir les données~?
    \item Appliquez la grille CRO à la décision d'invoquer l'art.~23
          ou de passer directement au MLAT.
\end{enumerate}
\end{boxexo}

\begin{boxexo}[title={\ding{45}\quad Exercice~11.3 --- Stratégie juridique complète \niveauexpert}]
Dans l'Opération WOURI (Chapitre~\ref{chap:gr-affaires}), l'analyse
révèle que l'attaque BEC implique~:
\begin{itemize}
    \item Email phishing envoyé depuis un compte Outlook~(\texttt{@hotmail.fr})
          hébergé en Irlande (centre de données Microsoft Europe)~;
    \item VPN utilisé~: NordVPN (siège social~: Panama, serveurs
          aux Pays-Bas)~;
    \item Compte bancaire destinataire~: HSBC Hong Kong~;
    \item Fonds convertis en BTC et envoyés vers une adresse dont la
          dernière transaction pointe vers un exchange turc (Paribu).
\end{itemize}
\begin{enumerate}[(a)]
    \item Pour chaque élément, identifiez le cadre légal applicable et
          l'instrument de coopération optimal.
    \item La double incrimination est-elle satisfaite dans chaque cas~?
          Justifiez en citant les textes de loi des deux pays.
    \item Rédigez le plan d'action chronologique complet~: dans quel
          ordre agir, par quel canal, dans quel délai, et quel résultat
          attendre à chaque étape.
    \item Quel élément de cette chaîne est le plus vulnérable à un
          refus de coopération~? Proposez une stratégie de contournement.
\end{enumerate}
\end{boxexo}

\bigskip

\begin{boxinfo}
\textbf{Transition.} Le Chapitre~\ref{chap:droit-cam} approfondit le
\textbf{droit camerounais du numérique}~: architecture complète de
la \loicam{2010/012} et de la \loicam{2024/017}, jurisprudence naissante,
rôle de l'ANTIC, articulation avec le Code pénal et le Code de procédure
pénale. Ce chapitre sera le plus important pour la rédaction des
procès-verbaux et la qualification des infractions devant les tribunaux
camerounais.
\end{boxinfo}
